\section{Docker Compose}

Quando uma aplicação envolve múltiplos containers --- bancos de dados, caches, filas de mensagens, entre outros serviços --- surge a questão de como orquestrá-los. Uma boa prática é que cada container faça uma única função e a faça bem.

Embora seja possível usar vários comandos \texttt{docker run} para iniciar múltiplos containers, isso exige lidar manualmente com networks e diversas flags de configuração. O \textbf{Docker Compose} resolve esse problema ao permitir definir todos os containers e suas configurações em um único arquivo YAML. Qualquer pessoa que clone o repositório pode subir toda a aplicação com um único comando.

\subsection{Dockerfile vs Compose file}

O \textbf{Dockerfile} contém instruções para construir uma \textit{container image}, enquanto o \textbf{Compose file} define os containers que serão executados e suas configurações. Frequentemente, o arquivo Compose referencia um Dockerfile para buildar a imagem usada em um serviço específico.

\subsection{Estrutura do arquivo compose.yaml}

O arquivo \texttt{compose.yaml} (ou \texttt{docker-compose.yml}) define os serviços, redes e volumes da aplicação. A estrutura principal é:

\begin{verbatim}
services:
  nome-do-servico:
    image: imagem:tag           # imagem a ser usada
    build: ./caminho            # OU buildar a partir de um Dockerfile
    ports:
      - "8080:80"               # mapeamento HOST:CONTAINER
    environment:
      - VARIAVEL=valor          # variáveis de ambiente
    volumes:
      - meus-dados:/app/data    # montar volume
      - ./local:/app/config     # bind mount
    networks:
      - minha-rede              # rede a ser utilizada
    depends_on:
      - outro-servico           # dependência de inicialização

volumes:
  meus-dados:                   # declaração de volumes nomeados

networks:
  minha-rede:                   # declaração de redes customizadas
\end{verbatim}

\subsubsection{Campos principais}

\begin{itemize}
    \item \texttt{services}: define cada container da aplicação. Cada serviço é um container.
    \item \texttt{image}: especifica a imagem Docker a ser usada.
    \item \texttt{build}: alternativa ao \texttt{image} --- builda a imagem a partir de um Dockerfile local.
    \item \texttt{ports}: mapeamento de portas no formato \texttt{HOST:CONTAINER}.
    \item \texttt{environment}: variáveis de ambiente passadas ao container.
    \item \texttt{volumes}: montagem de volumes ou bind mounts no container.
    \item \texttt{networks}: redes às quais o serviço será conectado.
    \item \texttt{depends\_on}: define a ordem de inicialização dos serviços (indica que um serviço depende de outro para subir).
\end{itemize}

\subsection{Comandos básicos}

\begin{verbatim}
docker compose up -d --build    # subir todos os serviços em background
docker compose down             # derrubar todos os serviços
docker compose ps               # listar serviços em execução
docker compose logs             # ver logs de todos os serviços
docker compose logs -f <serv.>  # seguir logs de um serviço específico
\end{verbatim}

\subsection{Exemplo básico}

\begin{verbatim}
git clone https://github.com/dockersamples/todo-list-app
cd todo-list-app
docker compose up -d --build
docker compose down
\end{verbatim}

O arquivo \texttt{compose.yaml} dentro do projeto define todos os serviços, suas imagens, portas, volumes, networks e demais configurações necessárias.
